% Options for packages loaded elsewhere
% Options for packages loaded elsewhere
\PassOptionsToPackage{unicode}{hyperref}
\PassOptionsToPackage{hyphens}{url}
\PassOptionsToPackage{dvipsnames,svgnames,x11names}{xcolor}
%
\documentclass[
  letterpaper,
  DIV=11,
  numbers=noendperiod]{scrartcl}
\usepackage{xcolor}
\usepackage{amsmath,amssymb}
\setcounter{secnumdepth}{-\maxdimen} % remove section numbering
\usepackage{iftex}
\ifPDFTeX
  \usepackage[T1]{fontenc}
  \usepackage[utf8]{inputenc}
  \usepackage{textcomp} % provide euro and other symbols
\else % if luatex or xetex
  \usepackage{unicode-math} % this also loads fontspec
  \defaultfontfeatures{Scale=MatchLowercase}
  \defaultfontfeatures[\rmfamily]{Ligatures=TeX,Scale=1}
\fi
\usepackage{lmodern}
\ifPDFTeX\else
  % xetex/luatex font selection
\fi
% Use upquote if available, for straight quotes in verbatim environments
\IfFileExists{upquote.sty}{\usepackage{upquote}}{}
\IfFileExists{microtype.sty}{% use microtype if available
  \usepackage[]{microtype}
  \UseMicrotypeSet[protrusion]{basicmath} % disable protrusion for tt fonts
}{}
\makeatletter
\@ifundefined{KOMAClassName}{% if non-KOMA class
  \IfFileExists{parskip.sty}{%
    \usepackage{parskip}
  }{% else
    \setlength{\parindent}{0pt}
    \setlength{\parskip}{6pt plus 2pt minus 1pt}}
}{% if KOMA class
  \KOMAoptions{parskip=half}}
\makeatother
% Make \paragraph and \subparagraph free-standing
\makeatletter
\ifx\paragraph\undefined\else
  \let\oldparagraph\paragraph
  \renewcommand{\paragraph}{
    \@ifstar
      \xxxParagraphStar
      \xxxParagraphNoStar
  }
  \newcommand{\xxxParagraphStar}[1]{\oldparagraph*{#1}\mbox{}}
  \newcommand{\xxxParagraphNoStar}[1]{\oldparagraph{#1}\mbox{}}
\fi
\ifx\subparagraph\undefined\else
  \let\oldsubparagraph\subparagraph
  \renewcommand{\subparagraph}{
    \@ifstar
      \xxxSubParagraphStar
      \xxxSubParagraphNoStar
  }
  \newcommand{\xxxSubParagraphStar}[1]{\oldsubparagraph*{#1}\mbox{}}
  \newcommand{\xxxSubParagraphNoStar}[1]{\oldsubparagraph{#1}\mbox{}}
\fi
\makeatother


\usepackage{longtable,booktabs,array}
\usepackage{calc} % for calculating minipage widths
% Correct order of tables after \paragraph or \subparagraph
\usepackage{etoolbox}
\makeatletter
\patchcmd\longtable{\par}{\if@noskipsec\mbox{}\fi\par}{}{}
\makeatother
% Allow footnotes in longtable head/foot
\IfFileExists{footnotehyper.sty}{\usepackage{footnotehyper}}{\usepackage{footnote}}
\makesavenoteenv{longtable}
\usepackage{graphicx}
\makeatletter
\newsavebox\pandoc@box
\newcommand*\pandocbounded[1]{% scales image to fit in text height/width
  \sbox\pandoc@box{#1}%
  \Gscale@div\@tempa{\textheight}{\dimexpr\ht\pandoc@box+\dp\pandoc@box\relax}%
  \Gscale@div\@tempb{\linewidth}{\wd\pandoc@box}%
  \ifdim\@tempb\p@<\@tempa\p@\let\@tempa\@tempb\fi% select the smaller of both
  \ifdim\@tempa\p@<\p@\scalebox{\@tempa}{\usebox\pandoc@box}%
  \else\usebox{\pandoc@box}%
  \fi%
}
% Set default figure placement to htbp
\def\fps@figure{htbp}
\makeatother





\setlength{\emergencystretch}{3em} % prevent overfull lines

\providecommand{\tightlist}{%
  \setlength{\itemsep}{0pt}\setlength{\parskip}{0pt}}



 


\KOMAoption{captions}{tableheading}
\makeatletter
\@ifpackageloaded{caption}{}{\usepackage{caption}}
\AtBeginDocument{%
\ifdefined\contentsname
  \renewcommand*\contentsname{Table of contents}
\else
  \newcommand\contentsname{Table of contents}
\fi
\ifdefined\listfigurename
  \renewcommand*\listfigurename{List of Figures}
\else
  \newcommand\listfigurename{List of Figures}
\fi
\ifdefined\listtablename
  \renewcommand*\listtablename{List of Tables}
\else
  \newcommand\listtablename{List of Tables}
\fi
\ifdefined\figurename
  \renewcommand*\figurename{Figure}
\else
  \newcommand\figurename{Figure}
\fi
\ifdefined\tablename
  \renewcommand*\tablename{Table}
\else
  \newcommand\tablename{Table}
\fi
}
\@ifpackageloaded{float}{}{\usepackage{float}}
\floatstyle{ruled}
\@ifundefined{c@chapter}{\newfloat{codelisting}{h}{lop}}{\newfloat{codelisting}{h}{lop}[chapter]}
\floatname{codelisting}{Listing}
\newcommand*\listoflistings{\listof{codelisting}{List of Listings}}
\makeatother
\makeatletter
\makeatother
\makeatletter
\@ifpackageloaded{caption}{}{\usepackage{caption}}
\@ifpackageloaded{subcaption}{}{\usepackage{subcaption}}
\makeatother
\usepackage{bookmark}
\IfFileExists{xurl.sty}{\usepackage{xurl}}{} % add URL line breaks if available
\urlstyle{same}
\hypersetup{
  pdftitle={Causal Inference Notation Guide},
  pdfauthor={Davood Tofighi, Ph.D.},
  colorlinks=true,
  linkcolor={blue},
  filecolor={Maroon},
  citecolor={Blue},
  urlcolor={Blue},
  pdfcreator={LaTeX via pandoc}}


\title{Causal Inference Notation Guide}
\author{Davood Tofighi, Ph.D.}
\date{}
\begin{document}
\maketitle


This document summarizes the standard notation used throughout the
Regression Analysis course.

\subsection{General Principles}\label{general-principles}

A critical convention in statistics is distinguishing between a
\emph{random variable} and its \emph{observed value}.

\begin{itemize}
\item
  \textbf{Random Variables}: These are represented by \textbf{capital
  letters} (e.g., \(Y\), \(X_j\)). They represent the full range of
  potential outcomes of a variable before data is collected. Estimators
  (like \(\hat{\beta}_j\)) are also random variables, as their value
  depends on the particular sample drawn.
\item
  \textbf{Observed Values (Realizations)}: These are represented by
  \textbf{lowercase letters} (e.g., \(y_i\), \(x_{ij}\)). They are the
  specific, fixed numerical values that have been collected in a sample.
\end{itemize}

This document provides a reference guide to the mathematical notations
used in this course. Consistency in notation is key to clearly defining
causal questions and understanding the assumptions needed to answer
them.

\begin{center}\rule{0.5\linewidth}{0.5pt}\end{center}

\begin{equation}\phantomsection\label{eq-wq}{
\Var(X)
}\end{equation}

\subsubsection{Core Variables}\label{core-variables}

These are the fundamental building blocks of our causal models. We use
single uppercase letters to represent \textbf{observed variables}---the
actual data we collect for each individual.

\begin{itemize}
\tightlist
\item
  \textbf{\(Y\)}: The \textbf{outcome} variable of interest (e.g.,
  health status, income).
\item
  \textbf{\(A\)}: The \textbf{treatment} or exposure variable (e.g.,
  taking a drug, receiving a scholarship). This is often binary (\(1\)
  for treated, \(0\) for control).
\item
  \textbf{\(Z\)}: An \textbf{instrumental variable}, which is a variable
  that influences the treatment but not the outcome directly (e.g., a
  doctor's encouragement, a draft lottery number).
\item
  \textbf{\(X\) or \(C\)}: A vector of measured \textbf{covariates} or
  confounders (e.g., age, sex, pre-existing conditions).
\item
  \textbf{\(U\)}: A vector of \textbf{unmeasured confounders}, which are
  variables that affect both \(A\) and \(Y\) but are not observed in the
  data.
\end{itemize}

\begin{center}\rule{0.5\linewidth}{0.5pt}\end{center}

\subsubsection{Potential Outcomes
Notation}\label{potential-outcomes-notation}

This notation, central to the Neyman-Rubin causal model, is used to
describe hypothetical ``what-if'' scenarios.

\begin{itemize}
\tightlist
\item
  \textbf{\(Y(a)\)}: Represents the \textbf{potential outcome} for an
  individual if their treatment level \emph{were set to} \(a\).

  \begin{itemize}
  \tightlist
  \item
    \textbf{\(Y(1)\)}: The potential outcome if the individual receives
    the treatment.
  \item
    \textbf{\(Y(0)\)}: The potential outcome if the individual does not
    receive the treatment (control). The causal effect for an individual
    is \(Y(1) - Y(0)\), which is fundamentally unobservable.
  \end{itemize}
\item
  \textbf{\(A(z)\)}: Represents the \textbf{potential treatment} an
  individual would choose if their instrument \emph{were set to} \(z\).

  \begin{itemize}
  \tightlist
  \item
    \textbf{\(A(1)\)}: The treatment status if the instrument is present
    (\(Z=1\)).
  \item
    \textbf{\(A(0)\)}: The treatment status if the instrument is absent
    (\(Z=0\)).
  \end{itemize}
\item
  \textbf{Consistency}: This crucial assumption links the unobservable
  potential outcomes to the observed data. It states that the outcome we
  see is the potential outcome corresponding to the treatment that was
  actually received.

  \begin{itemize}
  \tightlist
  \item
    \(Y = Y(A)\)
  \item
    \(A = A(Z)\)
  \end{itemize}
\end{itemize}

\begin{center}\rule{0.5\linewidth}{0.5pt}\end{center}

\subsubsection{Probability and
Expectation}\label{probability-and-expectation}

These are standard notations from statistics used to describe
distributions and averages.

\begin{longtable}[]{@{}
  >{\raggedright\arraybackslash}p{(\linewidth - 6\tabcolsep) * \real{0.2261}}
  >{\raggedright\arraybackslash}p{(\linewidth - 6\tabcolsep) * \real{0.1307}}
  >{\raggedright\arraybackslash}p{(\linewidth - 6\tabcolsep) * \real{0.0955}}
  >{\raggedright\arraybackslash}p{(\linewidth - 6\tabcolsep) * \real{0.5477}}@{}}
\toprule\noalign{}
\begin{minipage}[b]{\linewidth}\raggedright
Symbol
\end{minipage} & \begin{minipage}[b]{\linewidth}\raggedright
LaTeX Command
\end{minipage} & \begin{minipage}[b]{\linewidth}\raggedright
Name / Abbreviation
\end{minipage} & \begin{minipage}[b]{\linewidth}\raggedright
Description
\end{minipage} \\
\midrule\noalign{}
\endhead
\bottomrule\noalign{}
\endlastfoot
\(\textcolor{green}{\operatorname{E}[X]}\) &
\texttt{\textbackslash{}operatorname\{E\}{[}X{]}} & Expectation & The
expected value (or long-run average) of a random variable \(X\). \\
\(\textcolor{green}{\operatorname{Var}(X)}\) &
\texttt{\textbackslash{}operatorname\{Var\}(X)} & Variance & A measure
of the spread of a random variable \(X\). Defined as
\(\operatorname{E}[(X - \operatorname{E}[X])^2]\). \\
\(\textcolor{green}{\operatorname{Cov}(X, Y)}\) &
\texttt{\textbackslash{}operatorname\{Cov\}(X,\ Y)} & Covariance & A
measure of the joint variability of two random variables \(X\) and
\(Y\). \\
\(\textcolor{green}{\operatorname{SD}(X)}\) &
\texttt{\textbackslash{}operatorname\{SD\}(X)} & Standard Deviation &
The square root of the variance, \(\sqrt{\operatorname{Var}(X)}\),
expressed in the original units of \(X\). \\
\(\textcolor{green}{\operatorname{Cor}(X, Y)}\) &
\texttt{\textbackslash{}operatorname\{Cor\}(X,\ Y)} & Correlation & A
standardized measure of the linear relationship between two random
variables \(X\) and \(Y\). \\
\(\textcolor{green}{\operatorname{SE}(X)}\) &
\texttt{\textbackslash{}operatorname\{SE\}(X)} & Standard Error & A
measure of the variability of a statistic (e.g., a regression
coefficient) across different samples. \\
\end{longtable}

\begin{itemize}
\tightlist
\item
  \textbf{\(P(\cdot)\)}: The \textbf{probability} of an event.
\item
  \textbf{\(\operatorname{E}[\cdot|\cdot]\)}: The \textbf{conditional
  expectation}. For example, \(\operatorname{E}[Y|A=1]\) is read as
  ``the expected outcome among the subgroup of individuals who were
  observed to have received the treatment.'' This represents an
  \textbf{association}.
\item
  \textbf{\(\perp\!\!\!\perp\)}: Represents \textbf{statistical
  independence} between two variables. For example,
  \(Z \perp\!\!\!\perp U\) means the instrument is independent of
  unmeasured confounders.
\end{itemize}

\begin{center}\rule{0.5\linewidth}{0.5pt}\end{center}

\subsubsection{Causal Estimands}\label{causal-estimands}

These are the specific causal quantities we aim to estimate from our
data.

\begin{itemize}
\tightlist
\item
  \textbf{ATE (Average Treatment Effect)}: The average causal effect for
  the entire population. \[ATE = \operatorname{E}[Y(1) - Y(0)]\]
\item
  \textbf{LATE (Local Average Treatment Effect)}: The average causal
  effect specifically for the subgroup of ``compliers''---individuals
  whose treatment choice is changed by the instrument.
  \[LATE = \operatorname{E}[Y(1) - Y(0) | A(1) > A(0)]\]
\item
  \textbf{ITT (Intent-to-Treat Effect)}: The effect of the
  \emph{assignment} of the instrument on the outcome, regardless of
  whether the treatment was actually taken.
  \[ITT = \operatorname{E}[Y|Z=1] - \operatorname{E}[Y|Z=0]\]
\end{itemize}

\begin{center}\rule{0.5\linewidth}{0.5pt}\end{center}

\subsubsection{\texorpdfstring{The
\texttt{do}-operator}{The do-operator}}\label{the-do-operator}

This notation, from the work of Judea Pearl, provides an alternative way
to express causal interventions. It distinguishes ``seeing'' from
``doing.''

\begin{itemize}
\tightlist
\item
  \textbf{\(\operatorname{E}[Y|\text{do}(A=a)]\)}: Represents the
  expected outcome \(Y\) if we were to \textbf{intervene} and set the
  treatment \(A\) to the value \(a\) for everyone in the population.
  This represents a \textbf{causal} effect.
\item
  \textbf{Equivalence}: A \texttt{do}-expression is equivalent to a
  potential outcome expectation.
  \[\operatorname{E}[Y|\text{do}(A=a)] = \operatorname{E}[Y(a)]\]
\end{itemize}




\end{document}
